% !TeX root = ../main.tex
\chapter{Anforderungsanalyse}
\label{cha:anforderungen}
% Das wichtigste Kapitel für den roten Faden: Konzipierung, Implementierung,
% Tests und die abschließende Bewertung dürfen sich NUR auf die hier
% dokumentierten Anforderungen stützen. Anforderungen mit IDs versehen
% (FA-1, NFA-1, ...), damit sie sich in den Folgekapiteln zitieren lassen.

\section{Funktionale Anforderungen}
\label{sec:fa}

\autoref{tab:fa} listet die funktionalen Anforderungen.

\begin{table}[H]
    \centering
    \begin{tabular}{p{1.4cm}p{12cm}}
        \hline
        \textbf{ID} & \textbf{Anforderung} \\
        \hline
        FA-1 & \dots \\
        FA-2 & \dots \\
        \hline
    \end{tabular}
    \caption{Funktionale Anforderungen}
    \label{tab:fa}
\end{table}

\section{Nicht-funktionale Anforderungen}
\label{sec:nfa}
